% ============================================================================
% IX. CONCLUSION  (one paragraph, formal register)
% ============================================================================
\section{Conclusion}\label{sec:conclusion}

% ---- P-IX.1 -------------------------------------------------------------------
% PURPOSE: three sentences of findings and one of outlook.
% CLAIMS: (a) equivalence under A1--A3 explains the v1 tie, verified on the
%   real sweep (10,000/10,000 off-grid decisions; 11,300/11,300 for B9+)
%   [V1T]; (b) the boundary: C1 second order, C2 first order under commit,
%   belief-gated speculative prefetch closes the gap [TH]; (c) the v2
%   boundary test (P1, P2, P3)
%   outcome [docs/research_v2/v2_test_results_final.md]: on d1_open
%   (pre-registered per-dataset analysis) the boundary claim is not
%   supported -- P1 fails, P2 fails Holm, P3 fails -- and H-REV holds:
%   the decomposed B9+ beats the joint controller under commit,
%   Delta_C = -0.0367 [-0.0481, -0.0252], p < 1e-4; pooled with
%   d2_support (exploratory) the verdicts are the same (Delta_C -0.0351
%   [-0.0458, -0.0245]). CONTRADICTS (b): "belief-gated speculative
%   prefetch closes the gap" is theory only; v2 did not show it; (d) outlook: real recorded
%   speech, other generators, and learned end-of-turn prediction.
% NUMBERS: 10,000/10,000; 11,300/11,300 [V1T].
\todo{P-IX.1: conclusion.}

% ---- Acknowledgment (optional; IEEE spelling "Acknowledgment") ----------------
% \section*{Acknowledgment}
% \PEND{compute credits / funding, if any; no AI-tool acknowledgment wording
%   unless the conference requires a disclosure statement -- check the
%   INDICON author kit}
